\chapter{Universo matematico e controparte non standard}\stepcounter{capitoloeccolo}

% Dopo aver mostrato la costruzione di $ \starR $ come ultraprodotto, si utilizza una costruzione simile, generalizzandola ad un \emph{universo matematico}. Questo consentirà di superare parzialmente la limitazione della scelta del linguaggio $ \mathcal{L} $, e quindi di parlare degli oggetti matematici ``comuni'' e della loro controparte non standard, pur senza uscire dal formalismo logico. 

%TODO: Mostrare il perché voglio fare quest costruzione:
% Noto che, poiché l'obiettivo di studio di questa tesi sono i manifolds, allora non è sufficiente studiare semplicemente R*, bensì c'è bisogno di una struttura più grande. Si introdurrà quindi il concetto di sovrastruttura di un insieme X, e se ne costruirà *V(X). 
% Si riconoscerà poi un'immersione elementare tra R* (come definita nel chap precedente) e R* all'interno di *V(R). Vedi retro tesi
% All'interno di V(RR), poi, si identificherà il manifold M.

Nel capitolo precedente, a seguito della costruzione di $ \starR $, è stato necessario definire attentamente l'estensione di altri oggetti matematici strettamente collegati, come funzioni, relazioni d'ordine e sottoinsiemi.

Dunque la complessità nel definire tutti gli oggetti e le strutture necessarie per studiare le varietà differenziabili, ovvero l'obiettivo di questo lavoro, rende necessaria una generalizzazione della costruzione precedente tale per cui, preso un qualsiasi\footnote{Nella definizione formale si indicherà con più precisione cosa significhi questo ``qualsiasi''} oggetto matematico, sia chiaro e naturale quale sia la sua controparte non standard.

In particolare si costruirà una struttura $ V(\R) $ contentente tutti gli oggetti di interesse, e contenente $ \R $, e se ne definirà la controparte non standard $ \nonstandard{V(\R)} $. Si mostrerà che $ \starR $, come definito nel capitolo precedente, e con tutte le sue proprietà come campo ordinato, si immerge elementarmente in un elemento di $ \nonstandard{V(\R)} $. 

Questa è la ragione della apparente rindondanza del capitolo precedente: è più intuitivo svolgere le dimostrazioni in un contesto meno astratto, e solo infine dimostrare che i risultati continuano a valere cambiando leggermente le premesse iniziali.

\section{Universo matematico}

\definizione{
    Sia $ X $ un insieme generico. Siano \begin{equation*}
        V_0(X) = X, \qquad V_{n+1}(X) = V_{n}(X)\cup\parti{V_{n}(X) },\quad \forall\, n \in \N\setminus\{0\}.
    \end{equation*}Si definisce \emph{universo matematico} costruito su $ X $ l'insieme \begin{equation*}
        V(X)\coloneqq \bigcup_{n \in \N} V_{n}(X). 
    \end{equation*}
}
\osservazione{
    Valgono alcuni fatti la cui dimostrazione è immediata: \begin{itemize}
        \item se $ a \in V(X)\setminus X $, allora $ a \subseteq V(X) $;
        \item $ V_{m}(X) \subsetneqq V_{m+1}(X)$.
    \end{itemize}
}
\definizione{
    Un elemento $ x \in V(X) $ si dice \emph{atomico} se $ x\neq \emptyset $ e $ x\cap V(X) = \emptyset $.


    Un insieme $ X $ si dice \emph{insieme base} se $ X\neq \emptyset $ e ogni suo elemento è atomico. 
}
\paragrafo{Notazione}{%
    D'ora in avanti $ X $ indicherà sempre, salvo diversamente indicato, un insieme base.
}{}{}
\paragrafo{Corollario}{%
    Sia $ X $ è un insieme base. Se $ a \in V(X) $, $ a \in b \in V_{m}(X)  $, allora $ m\ge 1 $ e $ a \in V_{m-1}(X) $.
}{}{}

Il corollario precedente è il 4.4.1 di \cite{chang_model_2012}.

\paragrafo{Corollario}{%
    Sia $ X $ un insieme base. Se $ a,b \in V(X) $ e $ a\cap V(X) = b\cap V(X) $, allora necessariamente $ a=b $ oppure $ a, b \in X\cup\{\emptyset\}$. 
}{dkjncskjndjjdjskjnckjnsjdkkjnkjnkjnknkkjnkjnjnj}{}
\dimframmento{dkjncskjndjjdjskjnckjnsjdkkjnkjnkjnknkkjnkjnjnj}{
    Se $ a \notin X $, allora $ a \in \bigcup_{m\ge 1} V_{m}(X) $, dunque $ a \subseteq V(X) $, e quindi $ a \cap V(X) = a $.
    \begin{itemize}
        \item Se $ b \notin X $, si ha che $ b\cap V(X) = b $, come sopra, e dunque $ a=b $;
        \item Se $ b \in X $, allora $ b\cap V(X) = \emptyset $ (poiché $ X $ è un insieme base), e pertanto $ a = \emptyset $.\qed
    \end{itemize}
}
\osservazione{
    L'insieme dei numeri reali $ \R $ è un insieme base. Questo sarà estremamente importante ai fini di questa tesi, ma la dimostrazione prescinde dallo scopo del lavoro. Pertanto, si rimanda alla sezione 16.2 di \cite{katz_true_nodate}.
}
\section{Linguaggio del prim'ordine}

La costruzione di un universo matematico consente di quantificare sugli insiemi, pur mantentendosi nell'ambito della logica del prim'ordine. 

In particolare, ogni formula sarà espressa nel linguaggio \begin{equation*}
    \mathcal{L} = \{\in\}
\end{equation*}Per chiarezza, le $ \mathcal{L} $-formule verranno chiamate $ \in $-formule.

\definizione{
    Una $ \in $-formula $ \phi $ è detta \emph{limitata} se e solo se i quantificatori $ \forall\, x $ e $ \exists\, x $ appaiono in $ \phi $ \emph{solo} nelle forme seguenti: \begin{equation*}
        \forall\,x\, (x \in a \,\implies\, \cdots)
    \end{equation*}oppure \begin{equation*}
        \exists\,x\, (x \in a \,\land\, \cdots)
    \end{equation*}dove $ a $ è una variabile oppure una costante, e $ \phi $ può contenere altre variabili libere.
}
\osservazione{
    Le forme di prima saranno abbreviate come segue: \begin{align*}
        \forall\,x\, (x \in a \,\implies\, \cdots)\quad &\leadsto\quad\forall\, x \in a\, (\cdots)\\ 
        \exists\,x\, (x \in a \,\land\, \cdots) \quad &\leadsto \quad \exists\, x \in a\, (\cdots).
    \end{align*}
}
\section{Costruzione dell'universo non standard come ultraprodotto}

\definizione{
    Una estensione non standard di $ V(X) $ consiste in un insieme base $ Y $, con $ X \subsetneqq Y $, un sotto universo $ \nonstandard{V(X)} \subseteq V(Y) $, e una funzione \begin{equation*}
        *: V(X) \longrightarrow V(Y)
    \end{equation*}che soddisfi le seguenti condizioni. \begin{enumerate}
        \item[(I)] $ Y \subseteq \nonstandard{V(X)} $, e $ * $ è una mappa iniettiva tale per cui, se $ x \in X $, allora $ \nonstandard{x} \in Y $. 
        \item[(II)] Per ogni formula limitata $ \phi(x_1,\dots,x_{k} ) $ e ogni $ k $-upla di costanti $ a_1,\dots,a_{k} \in V(X)  $, la formula $ \phi[a_1,\dots,a_{k} ] $ è vera in $ V(X) $ se e solo se $ \phi[\nonstandard{a_1},\dots,\nonstandard{a_{k} }] $ è vera in $ \nonstandard{V(X)} $.
    \end{enumerate}
}
\osservazione{
    Si ha che $ Y = \nonstandard{X} $\footnote{Non riesco a convincermi di questa cosa, però non avrebbe senso il contrario}. %TODO capire se è vero
}

Lo scopo di tutta questa sezione sarà dimostrare che esistono estensioni non standard di $ V(X) $, tramite la costruzione dell'ultraprodotto ridotto.

\paragrafo{Notazione}{%
    D'ora in avanti, sia $ \mathcal{F} \subseteq \parti{\N}$ un arbitrario ultrafiltro di $ \N $, ed in particolare, un ultrafiltro \emph{non principale}. La sua esistenza è garantita dalla proposizione \framref{dkdkjncjdjdkjncjskjncjdkjncjdkjncjskjndkjckjkjnsd}.
}{}{}

\subsection{Ultraprodotto ridotto.}